\chapter{GİRİŞ}
Giriş bölümü, okuyucunun konuyu anlayıp projenin amacını ve konuya nasıl bir katkıda bulunduğunu değerlendirilmesi için yeterli temel bilgileri içermelidir. Bu amaçla projenin konusu tanımlanmalı, çalışmanın yapılmasının gereği, amacı ve hedefi kısaca anlatılmalıdır. Projenin motivasyonu; bir başka deyişle, bu konunun seçiliş sebebi ve konunun neden önemli olduğu, giriş bölümünde iyi bir şekilde vurgulanmalıdır. Ayrıca projenin konuya olan özgün katkısı varsa giriş bölümünde açık bir şekilde anlatılmalıdır. Bunların yanında proje çalışmasının anlaşılabilmesi için bilinmesi gereken ön bilgiler, giriş bölümünde okuyucuya aktarılmalıdır. Giriş bölümünün sonunda okuyucunun hangi bölümleri okuyacağına karar vermesini kolaylaştırmak için proje kitabının sonraki bölümleri kısaca tanıtılmalıdır.

\section{Amaç}

\section{Ön İnceleme}




